\documentclass[11pt]{article}

% Change "review" to "final" to generate the final (sometimes called camera-ready) version.
% Change to "preprint" to generate a non-anonymous version with page numbers.
\usepackage[review]{acl}
% Standard package includes
\usepackage{times}
\usepackage{latexsym}
\usepackage{booktabs}
\usepackage{amsmath}
\usepackage{multirow}
\usepackage{amssymb}

% For proper rendering and hyphenation of words containing Latin characters (including in bib files)
\usepackage[T1]{fontenc}
% For Vietnamese characters
% \usepackage[T5]{fontenc}
% See https://www.latex-project.org/help/documentation/encguide.pdf for other character sets

% This assumes your files are encoded as UTF8
\usepackage[utf8]{inputenc}

% This is not strictly necessary, and may be commented out,
% but it will improve the layout of the manuscript,
% and will typically save some space.
\usepackage{microtype}

% This is also not strictly necessary, and may be commented out.
% However, it will improve the aesthetics of text in
% the typewriter font.
\usepackage{inconsolata}

%Including images in your LaTeX document requires adding
%additional package(s)
\usepackage{graphicx}

% If the title and author information does not fit in the area allocated, uncomment the following
%
%\setlength\titlebox{<dim>}
%
% and set <dim> to something 5cm or larger.

\title{Beyond Redundancy: Diagnosing Information Distribution in Text Embeddings via Task-Aware Dimension Selection}

% Author information can be set in various styles:
% For several authors from the same institution:
% \author{Author 1 \and ... \and Author n \\
%         Address line \\ ... \\ Address line}
% if the names do not fit well on one line use
%         Author 1 \\ {\bf Author 2} \\ ... \\ {\bf Author n} \\
% For authors from different institutions:
% \author{Author 1 \\ Address line \\  ... \\ Address line
%         \And  ... \And
%         Author n \\ Address line \\ ... \\ Address line}
% To start a separate ``row'' of authors use \AND, as in
% \author{Author 1 \\ Address line \\  ... \\ Address line
%         \AND
%         Author 2 \\ Address line \\ ... \\ Address line \And
%         Author 3 \\ Address line \\ ... \\ Address line}

\author{First Author \\
  Affiliation / Address line 1 \\
  Affiliation / Address line 2 \\
  Affiliation / Address line 3 \\
  \texttt{email@domain} \\\And
  Second Author \\
  Affiliation / Address line 1 \\
  Affiliation / Address line 2 \\
  Affiliation / Address line 3 \\
  \texttt{email@domain} \\}

%\author{
%  \textbf{First Author\textsuperscript{1}},
%  \textbf{Second Author\textsuperscript{1,2}},
%  \textbf{Third T. Author\textsuperscript{1}},
%  \textbf{Fourth Author\textsuperscript{1}},
%\\
%  \textbf{Fifth Author\textsuperscript{1,2}},
%  \textbf{Sixth Author\textsuperscript{1}},
%  \textbf{Seventh Author\textsuperscript{1}},
%  \textbf{Eighth Author \textsuperscript{1,2,3,4}},
%\\
%  \textbf{Ninth Author\textsuperscript{1}},
%  \textbf{Tenth Author\textsuperscript{1}},
%  \textbf{Eleventh E. Author\textsuperscript{1,2,3,4,5}},
%  \textbf{Twelfth Author\textsuperscript{1}},
%\\
%  \textbf{Thirteenth Author\textsuperscript{3}},
%  \textbf{Fourteenth F. Author\textsuperscript{2,4}},
%  \textbf{Fifteenth Author\textsuperscript{1}},
%  \textbf{Sixteenth Author\textsuperscript{1}},
%\\
%  \textbf{Seventeenth S. Author\textsuperscript{4,5}},
%  \textbf{Eighteenth Author\textsuperscript{3,4}},
%  \textbf{Nineteenth N. Author\textsuperscript{2,5}},
%  \textbf{Twentieth Author\textsuperscript{1}}
%\\
%\\
%  \textsuperscript{1}Affiliation 1,
%  \textsuperscript{2}Affiliation 2,
%  \textsuperscript{3}Affiliation 3,
%  \textsuperscript{4}Affiliation 4,
%  \textsuperscript{5}Affiliation 5
%\\
%  \small{
%    \textbf{Correspondence:} \href{mailto:email@domain}{email@domain}
%  }
%}

\begin{document}
\maketitle

\begin{figure*}[!t]
    \centering
    \includegraphics[width=0.95\linewidth]{figures/fig_concept_overview.pdf}
    \caption{Research framework and main findings. We diagnose how task-relevant information is distributed across embedding dimensions by comparing optimized, random, and anti-optimized chunk selection within tasks, transferring chunk rankings across tasks, and analyzing chunk-importance concentration and entropy. Results show that task-relevant information is broadly distributed but not structureless: concentration varies by training paradigm, and task-specific rankings coexist with shared redundancy.    }
    \label{fig:concept_overview}
\end{figure*}

\begin{abstract}
Text embedding models often remain robust after substantial random dimension removal, suggesting considerable redundancy in high-dimensional representations. However, such robustness does not reveal how task-relevant information is distributed across dimensions: whether it is concentrated in a few critical dimensions or spread broadly across the embedding space. We diagnose this information distribution across 11 embedding models and 34 MTEB tasks by comparing random, task-aware optimized, and anti-optimized dimension selection, and by transferring dimension rankings across tasks. At 256 retained dimensions, optimized selection outperforms random by only +2.4--5.3\% for Retrieval-optimized Embedders and most Instruction-conditioned Embedders, but by +8.2--11.8\% for RoBERTa-InBedder and General-purpose Language Model Backbones. Under stronger compression, the optimized–random gap increases, indicating that broad redundancy coexists with exploitable structure across embedding dimensions. Cross-task transfer shows that dimension rankings are largely uncorrelated across tasks, yet transferred selections remain slightly above random across retained dimensionalities, revealing task-specific rankings together with shared redundancy. Entropy analysis shows near-uniform yet model-dependent importance distributions, explaining why many dimension subsets preserve performance while some models exhibit stronger concentration. Overall, task-aware dimension selection reveals a training-paradigm-dependent information distribution.
\end{abstract}

\section{Introduction}
Dense text embeddings are widely used in semantic search, clustering, classification, and reranking \citep{reimers-gurevych-2019-sentence, karpukhin-etal-2020-dense, khattab-zaharia-2020-colbert}, but their 768--4096-dimensional vectors impose substantial storage and computation costs at scale \citep{johnson-etal-2019-billion}. Recent work has shown that randomly removing half of the dimensions from modern text embeddings causes only modest degradation on retrieval and classification tasks \citep{takeshita2025randomly}, with similar robustness reported under more aggressive compression settings \citep{tsukagoshi2025redundancy}. These findings suggest considerable redundancy in text embedding representations.

However, redundancy alone does not reveal whether task-relevant information is concentrated in a small set of critical dimensions or broadly distributed across the embedding space. This distinction has a direct empirical implication: if information is broadly distributed, random dimension selection should perform close to task-aware optimized selection; if information is more concentrated, optimized selection should yield a clear advantage. We therefore use the optimized--random gap as a diagnostic of information distribution, and measure this gap across 11 embedding models and 34 MTEB tasks \citep{mteb} at different compression levels. 

Beyond asking how concentrated task-relevant information is within individual tasks, we also ask whether dimension importance is task-specific. If each task relies on its own set of important dimensions, then rankings derived from one task should transfer poorly to another. If tasks share redundant structure, transferred selections may preserve performance even when rankings disagree. We test this through cross-task transfer, ranking dimensions on a source task and evaluating the retained dimensions on a target task.

We observe that random selection already approaches task-aware selection under moderate compression for Retrieval-optimized Embedders and most Instruction-conditioned Embedders. For RoBERTa-InBedder and General-purpose Language Model Backbones, task-aware selection yields larger gains, suggesting stronger information concentration. Cross-task transfer further shows that dimension rankings vary substantially across tasks without causing proportional performance loss, indicating that shared redundancy buffers the impact of task-specific rankings. Figure~\ref{fig:concept_overview} summarizes our research framework and key findings.

Our contributions are:
\begin{enumerate}
\item We propose dimension selection as a diagnostic for studying how task-relevant information is distributed across embedding dimensions.
\item We show that task-relevant information distribution differs by training paradigm and is clearer under stronger compression.
\item We demonstrate that task-specific dimension rankings coexist with shared redundancy.
\end{enumerate}

\section{Related Work}
\subsection{Embedding Redundancy and Compressibility}

\paragraph{Evidence of redundancy.}
Recent work shows that text embeddings contain substantial dimensional redundancy. \citet{takeshita2025randomly} find that randomly removing 50\% of dimensions causes only modest degradation on retrieval and classification tasks. \citet{tsukagoshi2025redundancy} show that prompt-based LLM embeddings remain robust under aggressive dimensionality reduction. \citet{zhang2024evaluating} compare unsupervised dimensionality reduction methods, showing PCA can halve embedding size with limited loss.

\paragraph{Compression-oriented methods.} 
Several methods take advantage of the redundancy of embedding vectors. \citet{matryoshka2023} train embeddings where the first $k$ dimensions are optimized to be a good representation at any truncation level, with follow-up works adapting this idea to LLM embedding compression \citep{matryoshka-adaptor, smec2025}. Sparse-coding approaches further relax the contiguous-prefix constraint \citep{beyond-matryoshka2025}. 

\paragraph{Our distinction.}
These works show that embeddings are highly redundant and compressible, but treat compression as an objective. We instead use dimension selection as a diagnostic for studying whether task-relevant information is concentrated in a few critical dimensions or broadly distributed.

\begin{figure*}[t]
    \centering
    \includegraphics[width=0.95\linewidth]{figures/fig6_pruning_ratio_sweep.pdf}
    \caption{Retention across retained dimensionalities for three representative models. Optimized (green) retains the highest-ranked chunks, Anti-optimized (red) retains the lowest-ranked chunks, and Random (gray) provides the random baseline. The optimized--random gap remains modest for GTE-Large and Stella EN 400M under moderate compression, but is larger for RoBERTa-InBedder and widens under stronger compression. Shaded bands show 95\% bootstrap confidence intervals over tasks.}
    \label{fig:pruning_ratio_sweep}
\end{figure*}

\subsection{Task-specific Compression and Representation Structure}

\paragraph{Task-specific compression.}
Several methods try to exploit task-specific dimension importance. \citet{dime2025} use retrieval feedback to estimate dimension utility, and \citet{learning-to-select} predict dimension importance from query embeddings, while \citet{gaprune2025} use gradient alignment for domain-specific pruning.

\paragraph{Representation structure.}
The broader literature examines geometric and effective-dimensional properties of neural representations \citep{ethayarajh2019representation, li2018measuring, alemohammad2023intrinsic}, showing that high-dimensional representations often exhibit lower-dimensional or anisotropic structure across models and tasks.

\paragraph{Our distinction.}
Rather than only estimating which dimensions are important for one task, we ask whether dimension importance is task-specific, and whether different tasks still share redundant structure. We use cross-task transfer as a diagnostic, ranking dimensions on a source task and evaluating them on target tasks.

\section{Methodology}
\label{sec:method}
\subsection{Models and Tasks}
\paragraph{Models.}
We evaluate 11 text embedding models with diverse architectures and training paradigms, which we group into 3 families. \emph{General-purpose Language Model Backbones} include RoBERTa-Large \citep{roberta} and BART-Base \citep{bart}, which are not originally designed as embedding models. \emph{Instruction-conditioned Embedders} include RoBERTa-InBedder \citep{inbedder}, Qwen3-Embedding-0.6B \citep{qwen3-embedding}, Stella EN 400M \citep{stella}, and Instructor-Large \citep{instructor}, whose embedding ability is added or strengthened through post-training adaptation. \emph{Retrieval-optimized Embedders} include GTR-T5-Large \citep{gtr}, BGE-M3 \citep{bge}, GTE-Base, MxBai-Embed-Large, and GTE-Large \citep{gte}, which are trained directly for retrieval-style representation learning. Appendix~\ref{sec:coverage} summarizes which evaluations are available for each model.

\paragraph{Tasks.}
We use 34 benchmark tasks from MTEB \citep{mteb} covering 7 categories:
11 classification tasks (e.g., AmazonReviewsClassification, Banking77Classification),
3 clustering tasks (BiorxivClusteringS2S, MedrxivClusteringS2S, TwentyNewsgroupsClustering),
5 retrieval tasks (ArguAna, CQADupstackEnglishRetrieval, NFCorpus, SCIDOCS, SciFact),
3 reranking tasks (AskUbuntuDupQuestions, SciDocsRR, StackOverflowDupQuestions),
8 STS tasks (BIOSSES, SICK-R, STS12--STS16, STSBenchmark),
3 pair classification tasks (SprintDuplicateQuestions, TwitterSemEval2015, TwitterURLCorpus), and
1 summarization task (SummEval).

\subsection{Dimension Selection as a Diagnostic}

For a $D$-dimensional embedding, a reduced representation is defined by a retained dimension set $A\subseteq\{1,\ldots,D\}$ with $|A|=k$. For model $M$, task $T$, and retained dimensions $A$, let $S^{M,T}(A)$ denote the MTEB main metric (e.g., accuracy, nDCG@10) using only dimensions in $A$. We report retention relative to the full-dimensional representation:
\begin{equation}
R^{M,T}(A)=\frac{S^{M,T}(A)}{S^{M,T}(\mathrm{Full})}\times 100\%
\end{equation}
$R^{M,T}(A) > 100\%$ indicates that dimension selection may improve task performance relative to the full-dimensional baseline. Task averages are reported with 95\% bootstrap confidence intervals, and paired comparisons are conducted using paired $t$-tests and Wilcoxon signed-rank tests, with Holm--Bonferroni correction applied.

\begin{figure}[t]
    \centering
    \includegraphics[width=\linewidth]{figures/fig_opt_vs_rnd_scatter.pdf}
    \caption{Mean random retention (x-axis) versus mean optimized retention (y-axis) at 256 retained dimensions. The dashed identity line marks parity, and vertical distance above the line corresponds to optimized–random gap. Green marks General-purpose Language Model Backbones, gray marks Instruction-conditioned Embedders, and red marks Retrieval-optimized Embedders.}
    \label{fig:opt_random_combined_2}
\end{figure}

\paragraph{Information concentration.}
To test whether task-relevant information is concentrated in a few critical dimensions or broadly distributed, we compare random selection with task-aware optimized selection. Random selection samples $k$ dimensions uniformly without replacement. We repeat this $B=10$ times with different random seeds and report the mean retention:
\begin{equation}
\bar{R}_{\mathrm{rand}}^{M,T}(k)
= \frac{1}{B}\sum_{b=1}^{B}R^{M,T}(A_{\mathrm{rand}}^{(b)}(k))
\end{equation}
where $A_{\mathrm{rand}}^{(b)}(k)$ is the retained dimension set in the $b$-th trial.
For optimized selection, we divide the $D$ dimensions into $N=D/w$ contiguous chunks and retain $m=k/w$ chunks. We primarily use $w = 2$ to balance detail and cost, and further evaluate sensitivity to chunk size as an auxiliary validation. Each chunk is scored by keeping only that chunk:
\begin{equation}
q_i^{M,T}=S^{M,T}(C_i)
\end{equation}
where $C_i$ denotes the $i$-th chunk. We rank chunks by $q_i^{M,T}$ and retain the $m$ highest-scoring chunks:
\begin{equation}
A_{\mathrm{opt}}^{M,T}(k)=\operatorname{Top}_{m}\left(\{q_i^{M,T}\}_{i=1}^{N}\right)
\end{equation}
We define the optimized--random gap as:
\begin{equation}
\Delta_{\mathrm{gap}}^{M,T}(k)
=
R^{M,T}(A_{\mathrm{opt}}^{M,T}(k))
-
\bar{R}_{\mathrm{rand}}^{M,T}(k)
\end{equation}
where a smaller gap indicates that random subsets are nearly as informative as task-aware ones, indicating broadly distributed task-relevant information. A larger gap instead suggests stronger concentration in top-ranked chunks.


\paragraph{Lower bound.}
We also evaluate anti-optimized selection, retaining m lowest-scoring chunks:
\begin{equation}
A_{\mathrm{anti}}^{M,T}(k)=\operatorname{Bottom}_{m}\left(\{q_i^{M,T}\}_{i=1}^{N}\right)
\end{equation}
If anti-optimized selection substantially underperforms random selection, then dimensions do carry task-relevant information, although that information is broadly distributed enough for random selection to remain strong.

\subsection{Cross-task Transfer}
To test whether dimension importance is task-specific and whether different tasks nevertheless share redundant structure, we transfer dimension rankings across tasks. For source task $T_s$ and target task $T_t$, we rank chunks by source-task scores and evaluate the selected chunks on the target task:
\begin{equation}
R_{s\rightarrow t}^{M}(k)=R^{M,T_t}\left(\operatorname{Top}_{m}\left(\{q_i^{M,T_s}\}_{i=1}^{N}\right)\right)
\end{equation}

We compare each transfer result with random selection on the same target task:
\begin{equation}
\Delta_{s\rightarrow t}^{M}(k)
=
R_{s\rightarrow t}^{M}(k)
-
\bar{R}_{\mathrm{rand}}^{M,T_t}(k)
\end{equation}

We also compute source--target rank agreement:
\begin{equation}
\rho_{s,t}^{M}=\mathrm{Spearman}\left(\{q_i^{M,T_s}\}_{i=1}^{N},\{q_i^{M,T_t}\}_{i=1}^{N}\right)
\end{equation}

Together, $\Delta_{s\rightarrow t}^{M}(k)$ and $\rho_{s,t}^{M}$ disentangle functional transfer from ranking agreement. The former measures how well a source-derived ranking transfers to the target task relative to random, while the latter captures the rank correlation between source and target importance scores.

\subsection{Auxiliary Experiments}
We conduct auxiliary experiments to test whether the main diagnostic conclusions are robust under alternative compression settings, data splits, and importance definitions. The pruning-ratio sweep examines how the optimized--random gap and cross-task transfer behavior change as fewer dimensions are retained. Train/test split validation tests whether optimized rankings generalize beyond the data used to construct them. Mechanism analysis further uses entropy and cumulative importance to characterize whether dimension importance is near-uniform or concentrated.

We also test robustness to chunk-scoring granularity, random subset variability, practical selection baselines, and alternative ways of estimating dimension importance. 

\begin{figure*}[t]
    \centering
    \includegraphics[width=0.95\linewidth]{figures/fig5_transfer_and_rho_combined.pdf}
    \caption{Category-level cross-task transfer and rank agreement at 256 retained dimensions, excluding same-task source--target pairs. Rows denote source task categories and columns denote target task categories. N/A appears because the Summarization category contains only one task. (a) Transfer retention averaged over source--target task categories. Retention remains stable across category pairs. (b) Spearman rank correlations between source-task and target-task chunk-importance rankings, averaged over category pairs. Correlations are close to zero across most pairs. Together, these results show that task-specific rankings coexist with shared redundancy across tasks.
}
    \label{fig:transfer_and_rho}
\end{figure*}


\section{Experimental Results and Analysis}
\subsection{Training Paradigms Exhibit Different Concentration Profiles}
We first compare optimized, random and anti-optimized selection across compression levels to diagnose how concentrated task-relevant information is across embedding dimensions. Figure~\ref{fig:pruning_ratio_sweep} shows retention from a single retained dimension up to 75\% retention (768 of 1024 dimensions) for three representative models, covering moderate to aggressive compression regimes. For GTE-Large at 256 retained dimensions, Table~\ref{tab:category_summary} reports category-level retention, while Appendix~\ref{sec:per_task_gte} provides the full task-level results across all 34 tasks. 


\begin{table}[t]
    \centering
    \caption{Optimized–random retention gap at 256 retained dimensions. $\Delta$ reports the mean retention difference between optimized and random selection for each model with 95\% bootstrap CI. $d$ is the Cohen's $d$ for paired per-task differences. $p_{\text{holm}}$ reports Holm-Bonferroni corrected $p$ values across 11 tests.}
    \label{tab:optimized_vs_random}
    \small
    \resizebox{\columnwidth}{!}{%
    \begin{tabular}{lcccc}
        \toprule
        Model & $\Delta$ & 95\% CI & $d$ & $p_{\text{holm}}$ \\
        \midrule
        \multicolumn{5}{l}{\textit{Retrieval-optimized Embedders}} \\
        GTR-T5-Large        & +4.69\% & [+2.91\%, +6.61\%] & 0.84 & $<$0.001 \\
        BGE-M3              & +4.61\% & [+3.17\%, +6.07\%] & 1.04 & $<$0.001 \\
        GTE-Base            & +3.49\% & [+2.15\%, +4.97\%] & 0.82 & $<$0.001 \\
        MxBai-Embed-Large   & +2.59\% & [+1.51\%, +3.83\%] & 0.83 & $<$0.001 \\
        GTE-Large           & +2.41\% & [+1.48\%, +3.40\%] & 0.83 & $<$0.001 \\
        \midrule
        \multicolumn{5}{l}{\textit{Instruction-conditioned Embedders}} \\
        RoBERTa-InBedder    & +8.19\% & [+5.96\%, +10.62\%] & 1.17 & $<$0.001 \\
        Qwen3-Embedding-0.6B     & +5.30\% & [+3.54\%, +7.28\%]  & 1.00 & $<$0.001 \\
        Stella EN 400M      & +3.22\% & [+2.16\%, +4.44\%]  & 0.93 & $<$0.001 \\
        Instructor-Large    & +2.89\% & [+1.80\%, +4.06\%]  & 0.83 & $<$0.001 \\
        \midrule
        \multicolumn{5}{l}{\textit{General-purpose Language Model Backbones}} \\
        RoBERTa-Large       & +11.75\% & [+8.61\%, +15.53\%] & 1.16 & $<$0.001 \\
        BART-Base           & +10.08\% & [+6.30\%, +13.72\%] & 0.90 & $<$0.001 \\
        \bottomrule
    \end{tabular}%
    }
\end{table}

\paragraph{Information concentration differs by training paradigm.}
Optimized selection consistently outperforms random selection across all 11 models, but the optimized--random gap differs sharply by model family at 256 retained dimensions, as shown in Figure~\ref{fig:opt_random_combined_2} and Table~\ref{tab:optimized_vs_random}. The gaps are smallest among Retrieval-optimized Embedders ($+2.4\%$--$4.8\%$), with random selection retaining 95--98\% of full-dimensional performance. Instruction-conditioned Embedders show a similar pattern except for RoBERTa-InBedder: their gaps remain modest at $+2.9\%$--$5.3\%$, with random retention of 95--97\%. These small gaps therefore provide diagnostic evidence that task-relevant information is broadly distributed across dimensions for these models. In contrast, RoBERTa-InBedder and General-purpose Language Model Backbones show larger optimized--random gaps that would materially affect downstream performance ($+$8.2\%--11.8\%). For these models, task-relevant information is more concentrated in top-ranked chunks, suggesting that dimension concentration is closely associated with training paradigm. 

\begin{table}[t]
    \centering
    \caption{Mean retention (\%) by task category for GTE-Large at 256 retained dimensions. Random, Optimized, and Anti-optimized columns report retention ratios for the three strategies relative to the full-dimensional baseline, averaged over tasks within each category. Gap is Optimized--Random gap.}
    \label{tab:category_summary}
    \resizebox{\columnwidth}{!}{
    \begin{tabular}{lcccc}
        \toprule
        Task Category & Random & Optimized & Gap & Anti-optimized \\
        \midrule
        Classification & 97.3 & 101.4 & +4.1 & 84.2 \\
        Clustering & 95.1 & 99.4 & +4.3 & 81.5 \\
        Retrieval & 94.3 & 94.3 & +0.0 & 93.3 \\
        Reranking & 98.8 & 100.4 & +1.6 & 96.0 \\
        STS & 99.1 & 101.1 & +2.0 & 95.5 \\
        Pair Classification & 99.4 & 100.9 & +1.5 & 97.2 \\
        Summarization & 99.8 & 98.7 & -1.1 & 98.8 \\
        \bottomrule
    \end{tabular}
}
\end{table}


\paragraph{Smaller retained subsets increase sensitivity to chunk choice.}
The optimized--random gap widens as fewer dimensions are retained. At 256 retained dimensions, the gap is only $+2.4$--$5.3\%$ for Retrieval-optimized Embedders and Instruction-conditioned Embedders other than RoBERTa-InBedder. At 64 retained dimensions, the gap widens to $+5.1$--$9.4\%$ for the same model group. This pattern suggests that performance becomes more sensitive to chunk choice as the retained subset shrinks, making non-uniform dimension importance more visible under stronger compression. The transition point differs across model families, as detailed in Appendix~\ref{sec:gap_by_dim}.

\paragraph{Anti-optimized selection exposes weakly informative dimensions.}
Anti-optimized selection provides a complementary view of the same information-distribution question by probing the least informative dimensions. At 256 retained dimensions, it degrades retention to 84--91\% across the tested models, well below random selection. This shows that dimensions do carry task-relevant information. For Retrieval-optimized Embedders and most Instruction-conditioned Embedders, the small optimized--random gap indicates broadly distributed task-relevant information, but the much lower anti-optimized retention shows that bottom-ranked chunks contribute less to task-relevant information. For RoBERTa-InBedder and General-purpose Language Model Backbones, the larger optimized--random gap indicates stronger concentration in high-ranked chunks, making anti-optimized selection more damaging.


\subsection{Task-specific Rankings Coexist with Shared Redundancy}
\label{sec:cross_task_transfer}
The optimized--random gap diagnoses how task-relevant information is distributed within each task. Cross-task transfer asks a complementary question: whether dimension-importance rankings are task-specific, and whether selections derived from one task still preserve information for another. We therefore rank chunks on a source task, retain the top-ranked chunks, and evaluate the resulting subset on a target task. 


\paragraph{Dimension rankings are task-specific.}
To quantify ranking agreement, we compute Spearman correlations between source-task and target-task standalone chunk-score vectors over the same $N$ chunks. Figure~\ref{fig:transfer_and_rho}(b) reports the results, averaged within each source--target task-category pair. The correlations are close to zero across most category pairs, with low ranking agreement even for within-category transfer, indicating that source-task rankings provide little information about the relative ordering of chunks for target tasks. Dimension importance is therefore strongly task-specific.

%3行

\paragraph{Transferred performance reveals shared redundancy.}
Cross-task transfer consistently exceeds random selection across retained dimensionalities, as shown in Table~\ref{tab:cross_task_multidim}. The transfer gap is positive at every retained dimensionality, ranging from $+$0.7--1.6\% after excluding same-task source--target pairs, with Holm-corrected $p<0.001$ throughout. Appendix~\ref{sec:multidim_appendix} reports per-model breakdowns. Given the near-zero rank agreement above, tasks largely disagree on how dimensions should be ranked. Nevertheless, the consistently above-random transfer shows that source-derived selections can still preserve information useful for other tasks. This indicates that shared redundancy buffers the impact of task-specific rankings. 

\begin{table}[t]
    \centering
    \caption{Pooled cross-task transfer vs.\ random selection across 11 models. Transfer consistently exceeds random selection across retained dimensionalities. All gap includes diagonal self-transfer, while off-diagonal gap excludes same-task source--target pairs. $p_{\mathrm{holm}}$: Holm-corrected $p$-values for off-diagonal gaps. Per-model and category breakdowns are shown in Appendix~\ref{sec:multidim_appendix}.}
    \label{tab:cross_task_multidim}
    \small
    \resizebox{\columnwidth}{!}{%
    \begin{tabular}{lccccc}
        \toprule
        Retained Dims & Transfer & Random & All gap & Off-diagonal gap & $p_{\mathrm{holm}}$ \\
        \midrule
        16 & 64.8\% & 63.8\% & +1.0\% & +0.7\% & $<$0.001 \\
        32 & 76.3\% & 75.3\% & +1.0\% & +0.7\% & $<$0.001 \\
        64 & 86.7\% & 85.6\% & +1.0\% & +0.8\% & $<$0.001 \\
        128 & 94.3\% & 93.4\% & +0.9\% & +0.7\% & $<$0.001 \\
        256 & 99.1\% & 98.3\% & +0.8\% & +0.7\% & $<$0.001 \\
        512 & 101.5\% & 99.8\% & +1.7\% & +1.6\% & $<$0.001 \\
        \bottomrule
    \end{tabular}%
    }
\end{table}

\paragraph{The shared-redundancy pattern is not category-specific.}
Figure~\ref{fig:transfer_and_rho}(a) examines whether transfer retention is concentrated in specific source--target category pairs. The heatmap shows no clear diagonal advantage. Instead, retention varies more by target category than by source category, suggesting that category-level differences mainly reflect target-side difficulty rather than special source--target compatibility. Thus, shared redundancy persists across task categories.



\begin{figure*}[!t]
    \centering
    \includegraphics[width=0.95\linewidth]{figures/fig8_evidence_summary.pdf}
    \caption{(a) Top-$k$ concentration curves for 9 models. (b) Boxplots of normalized Shannon entropy for the same 9 models. Although all models have high entropy, RoBERTa-InBedder and General-purpose backbones show relatively higher top-$k$ concentration and lower entropy, indicating stronger concentration.}
    \label{fig:mechanism_summary}
\end{figure*}



\begin{table}[t]
    \centering
    \small
    \caption{Held-out validation of optimized selection. 
Chunks are ranked on 70\% of each task and evaluated on the remaining 30\%. 
\textbf{Holdout Gap} is optimized--random gap on held-out data; 
\textbf{Same-Data Gap} is the drop from same-data to held-out evaluation.}
    \label{tab:train_test_split}
    \resizebox{\columnwidth}{!}{%
    \begin{tabular}{@{}lccc@{}}
    \toprule
    \textbf{Model} & \textbf{Holdout Gap} & \textbf{Win\%} & \textbf{Same-Data Gap} \\
    \midrule
    \multicolumn{4}{l}{\textit{Retrieval-optimized Embedders}} \\
    GTE-Large & 1.74\% & 74.5\% & 0.67\% \\
    BGE-M3 & 2.70\% & 79.9\% & 1.91\% \\
    MxBai-Embed-Large & 1.69\% & 75.0\% & 0.90\% \\
    \midrule
    \multicolumn{4}{l}{\textit{Instruction-conditioned Embedders}} \\
    RoBERTa-InBedder & 4.02\% & 81.3\% & 4.17\% \\
    Qwen3-Embedding-0.6B & 2.42\% & 77.3\% & 2.88\% \\
    Stella EN 400M & 2.21\% & 80.5\% & 1.01\% \\
    Instructor-Large & 1.52\% & 71.4\% & 1.37\% \\
    \midrule
    \multicolumn{4}{l}{\textit{General-purpose Language Model Backbones}} \\
    RoBERTa-Large & 5.15\% & 87.2\% & 6.60\% \\
    BART-Base & 4.26\% & 81.8\% & 5.82\% \\
    \bottomrule
    \end{tabular}%
    }
\end{table}

\subsection{Auxiliary Robustness Checks}
\label{sec:aux_results}
\paragraph{Held-out validation controls for same-data ranking bias.}
A potential concern is that optimized selection may overstate dimension concentration because chunks are ranked and evaluated on the same data. We therefore conduct a train/test split validation, ranking chunks on a 70\% split and evaluating them on the held-out 30\%. As shown in Table~\ref{tab:train_test_split}, the same training-paradigm pattern persists: Retrieval-optimized Embedders and most Instruction-conditioned Embedders continue to show smaller holdout gaps ($+1.52\%$--$2.70\%$), whereas RoBERTa-InBedder and General-purpose Language Model Backbones show larger ones ($+4.02\%$--$5.15\%$). The held-out advantage is also consistent across tasks: optimized selection beats random in 71.4--87.2\% of held-out comparisons. These results support our diagnostic interpretation: the paradigm-dependent concentration pattern persists beyond the data used for ranking.

\paragraph{Additional robustness checks.}
We include several additional checks in the appendices. Appendix~\ref{sec:chunk_size} varies the chunk size used for scoring dimensions, testing whether the diagnosed concentration pattern is sensitive to chunk granularity. 
Appendices~\ref{sec:random_variance} and~\ref{sec:task_low_retention} quantify random-seed variability and task-level lower-tail outcomes, showing that moderate-compression conclusions are not driven by unstable seeds or hidden task-level failures.
Appendix~\ref{sec:practical_baselines} compares practical selection methods such as sequential, magnitude, gradient saliency, activation variance, learned masks, as well as projection-based compression baselines including PCA and random projection. These experiments verify the robustness of our main conclusions.


\section{Mechanism}
To further understand why the optimized--random gap differs across model families, we analyze the distribution of
standalone chunk importance through normalized Shannon entropy. For each model-task pair, the standalone importance of chunk $i$ is its task-specific evaluation score $q_i^{M,T}$, and $ p_i^{M,T} = q_i^{M,T}/\sum_j q_j^{M,T}$ is the normalized importance. We then compute:
\begin{equation}
H_{\text{norm}} = -\frac{1}{\log_2 N} \sum_{i=1}^{N} p_i^{M,T} \log_2 p_i^{M,T}
\end{equation}
where $N$ is the total number of chunks. Higher entropy indicates that importance is spread more evenly across chunks. Appendix~\ref{sec:loo} repeats the analysis with alternative importance definitions beyond standalone chunk scoring. Appendix~\ref{theoretical_perspective} provides a basis-dependence analysis, showing the same pattern appears under random orthogonal rotations, PCA, and whitening.

\paragraph{Top-$k$ concentration and entropy track model-family gaps.}
Retrieval-optimized embedders and most instruction-conditioned embedders exhibit flatter top-$k$ concentration curves and higher entropy, indicating that task-relevant information is distributed broadly. This pattern is consistent with their smaller optimized--random gaps, since random subsets are likely to retain enough useful signal when importance is near-uniform. By contrast, RoBERTa-InBedder and General-purpose Language Model Backbones consistently show relatively higher top-$k$ concentration curves and lower entropy, suggesting a more concentrated importance distribution and explaining why optimized selection provides larger gains for these models.

\paragraph{High-entropy rankings support transferability.}
The same mechanism also helps explain why cross-task transfer remains consistently above random despite low rank agreement. When many chunks have similar standalone importance, task-specific differences can substantially reorder the chunk rankings, producing low Spearman correlations across tasks. However, source-task rankings can still give higher priority to chunks that carry information useful across tasks. This produces a small but stable advantage over random selection. 

\section{Discussion}
\paragraph{Dimension selection reveals structured information distribution.}
Dimension selection provides a unified diagnostic for how task-relevant information is distributed within tasks and shared across tasks. Within a task, the optimized--random gap measures the degree of information concentration: smaller gaps indicate broadly distributed information, whereas larger gaps indicate stronger reliance on top-ranked chunks. Across tasks, weak rank agreement indicates that exact chunk-importance rankings are task-specific, while consistently above-random transfer gaps show shared redundancy in the embedding space. Rather than sharing a universal chunk ordering, tasks appear to share a coarser redundant structure: many chunks carry transferable task-relevant signal, while some chunks are repeatedly assigned low standalone importance. Together, these diagnostics give redundancy a more precise interpretation.

\paragraph{Training paradigms as a source of information distribution differences.}
\label{sec:training_paradigm_discussion}

Prior work shows that pretrained contextual representations and BERT sentence embeddings often exhibit anisotropic, non-uniform geometry, and that fine-tuning can further reshape embedding-space structure \citep{ethayarajh2019representation, li2020sentence, rajaee2021fine}. This helps explain why General-purpose Language Model Backbones, not trained specifically for similarity or retrieval, show larger optimized--random gaps: task-relevant information may follow a more uneven geometry, making top-ranked dimensions more informative. By contrast, Retrieval-optimized Embedders are trained directly for similarity-based representation learning, often with contrastive learning that affects alignment and uniformity \citep{wang2020understanding, gao2021representation}. This is consistent with their smaller optimized--random gaps, suggesting broader distribution of task-relevant information. RoBERTa-InBedder is a useful boundary case, with its stronger dimension structure potentially reflecting its short-answer-conditioned fine-tuning (Appendix~\ref{sec:Inbedder-mechanism}).

\section{Conclusion}

Prior work has shown that text embeddings are highly compressible, but compressibility alone does not reveal the representational structure behind this robustness. We use task-aware dimension selection to diagnose whether task-relevant information is broadly distributed or concentrated, how this pattern varies across training paradigms, and whether cross-task transfer reveals task-specific dimension importance or shared redundant structure.

Across 11 embedding models and 34 MTEB tasks, we find that Retrieval-optimized Embedders and most Instruction-conditioned Embedders show small optimized--random gaps, higher entropy, and flatter top-$k$ concentration curves, indicating broadly distributed task-relevant information. RoBERTa-InBedder and General-purpose Language Model Backbones show stronger concentration, suggesting that information distribution is closely associated with training paradigm. Cross-task transfer further shows that chunk-importance rankings are task-specific, yet transferred selections remain above random, indicating shared redundancy beyond exact rank agreement. Overall, dimension selection reveals a training-paradigm-dependent information distribution. Code and results will be released upon publication.

\label{sec:bibtex}

\section*{Limitations}
Our study provides a broad diagnostic analysis of task-relevant information distribution across 11 embedding models and 34 MTEB tasks. This coverage enables us to identify consistent patterns across model families, task categories, and compression levels, but it does not cover every possible deployment setting.

We focus on frozen, off-the-shelf text embedding models because this setting isolates how information is already organized in existing representations. Future work could extend the same diagnostic framework to models trained explicitly for adaptive dimensionality.

Our experiments study text embedding models. The same question may also be relevant for multi-modal embeddings, where representation geometry may differ. 

The observed association with training paradigm is empirical rather than theoretical. We do not provide a formal proof of this.

% Bibliography entries for the entire Anthology, followed by custom entries
%\bibliography{anthology,custom}
% Custom bibliography entries only
\bibliography{custom}
\clearpage

\appendix

\appendix
\section{Evaluation Coverage by Model}
\label{sec:coverage}
Table~\ref{tab:model_list} reports the output dimensionality of each model and summarizes which analyses are available for each of the 11 models in the study. For Stella EN 400M, which provides linear heads for multiple output sizes (512--8192), we use the default 1024-dimensional head before post-hoc dimension selection.

\begin{table*}[t]
    \centering
    \caption{Model coverage by experiment. ``Gap'' = optimized--random gap comparison and anti-optimized experiment. ``Transfer'' = cross-task transfer. ``Sweep'' = pruning-ratio sweep. ``Mechanism'' = mechanism analysis on entropy. ``Chunk'' = chunk size sweep analysis. ``Split'' = train/test split experiment. ``Tail'' = random-variance and tail-risk analysis. ``Magnitude'' = magnitude pruning evaluation. ``More Baselines'' = Practical baselines including Gradient Saliency, Activation Variance, PCA, and Learned Mask. ``More Definitions'' = alternative importance definitions analysis. ``Cost'' = retrieval cost analysis.}
    \label{tab:model_list}
    \scriptsize
    \resizebox{\textwidth}{!}{%
    \begin{tabular}{lcccccccccccc}
        \toprule
        Model & Dims & Gap & Transfer & Sweep & Mechanism & Chunk & Split & Tail & Magnitude & More Baselines & More Definitions & Cost \\
        \midrule
        RoBERTa-Large & 1024 & \checkmark & \checkmark & \checkmark & \checkmark &  & \checkmark & \checkmark &  & \checkmark & \checkmark & \checkmark \\
        BART-Base & 768 & \checkmark & \checkmark & \checkmark & \checkmark &  & \checkmark & \checkmark & \checkmark &  &  &  \\
        RoBERTa-InBedder & 1024 & \checkmark & \checkmark & \checkmark & \checkmark & \checkmark & \checkmark & \checkmark & \checkmark & \checkmark & \checkmark & \checkmark \\
        Qwen3-Embedding-0.6B & 1024 & \checkmark & \checkmark & \checkmark & \checkmark &  & \checkmark & \checkmark & \checkmark &  &  &  \\
        Stella EN 400M & 1024 & \checkmark & \checkmark & \checkmark & \checkmark & \checkmark & \checkmark & \checkmark & \checkmark & \checkmark & \checkmark & \checkmark \\
        Instructor-Large & 768 & \checkmark & \checkmark & \checkmark & \checkmark &  & \checkmark & \checkmark &  &  &  &  \\
        GTR-T5-Large & 768 & \checkmark & \checkmark & \checkmark &  & \checkmark &  &  &  &  &  &  \\
        BGE-M3 & 1024 & \checkmark & \checkmark & \checkmark & \checkmark &  & \checkmark & \checkmark & \checkmark & \checkmark & \checkmark & \checkmark \\
        GTE-Base & 768 & \checkmark & \checkmark & \checkmark &  & \checkmark &  &  & \checkmark &  &  &  \\
        MxBai-Embed-Large & 1024 & \checkmark & \checkmark & \checkmark & \checkmark &  & \checkmark & \checkmark & \checkmark &  &  &  \\
        GTE-Large & 1024 & \checkmark & \checkmark & \checkmark & \checkmark & \checkmark & \checkmark & \checkmark & \checkmark & \checkmark & \checkmark & \checkmark \\
        \bottomrule
    \end{tabular}%
    }
\end{table*}

\section{Per-Task Results for GTE-Large}
\label{sec:per_task_gte}
Table~\ref{tab:per_task_gte} reports per-task retention at 256 dimensions (corresponding to a 75\% dimensionality reduction) for all 34 tasks in our benchmark on GTE-Large. The source average asks whether a task provides a generally useful dimension ranking for other tasks, while the target average asks how well a task can be served by rankings derived from other tasks. Figure~\ref{fig:transfer_gte} further shows the corresponding cross-task transfer matrix. For every source-target task pair, using the source task's importance ranking retains 93--100\% of the target task's performance. These detailed results complement the main-text averages by showing how retention varies across individual tasks for each pruning method.

\begin{table*}[t]
    \centering
    \caption{Per-task retention (\%) at 256 retained dimensions for GTE-Large. Optimized and Anti-optimized use chunk size $w = 2$. Source average is the mean retention when the row task supplies the dimension ranking and all 34 tasks are targets. Target average is the mean retention when all 34 tasks supply rankings and the row task is the target.}
    \label{tab:per_task_gte}
    \small
    \resizebox{\textwidth}{!}{%
    \begin{tabular}{@{}p{5cm}ccccc@{}}
        \toprule
        Task & Random & Optimized & Anti-optimized & Source average & Target average  \\
        \midrule
        \multicolumn{6}{l}{\textit{Classification}} \\
        AmazonCounterfactualClassification & 95.8 & 105.6 & 70.0 & 97.9 & 96.5 \\
        AmazonReviewsClassification & 95.0 & 103.9 & 69.4 & 97.8 & 95.5 \\
        Banking77Classification & 99.3 & 100.0 & 98.1 & 97.5 & 99.3 \\
        EmotionClassification & 94.4 & 100.9 & 82.0 & 97.9 & 94.8 \\
        ImdbClassification & 99.3 & 100.9 & 77.5 & 97.1 & 99.4 \\
        MTOPDomainClassification & 99.0 & 99.9 & 95.5 & 98.0 & 99.0 \\
        MTOPIntentClassification & 97.8 & 99.2 & 94.7 & 97.5 & 97.9 \\
        MassiveIntentClassification & 98.0 & 99.5 & 96.7 & 97.8 & 98.1 \\
        MassiveScenarioClassification & 97.4 & 99.8 & 91.2 & 97.6 & 98.0 \\
        ToxicConversationsClassification & 96.9 & 105.0 & 66.6 & 97.5 & 97.3 \\
        TweetSentimentExtractionClassification & 97.9 & 101.2 & 84.8 & 97.7 & 97.8 \\
        \midrule
        \multicolumn{6}{l}{\textit{Clustering}} \\
        BiorxivClusteringS2S & 96.1 & 99.9 & 85.0 & 97.5 & 96.6 \\
        MedrxivClusteringS2S & 98.5 & 101.0 & 95.6 & 97.0 & 98.3 \\
        TwentyNewsgroupsClustering & 90.8 & 97.3 & 63.9 & 97.9 & 91.4 \\
        \midrule
        \multicolumn{6}{l}{\textit{Retrieval}} \\
        ArguAna & 96.7 & 96.7 & 96.7 & 97.4 & 96.8 \\
        CQADupstackEnglishRetrieval & 93.8 & 90.5 & 94.4 & 97.1 & 93.8 \\
        NFCorpus & 91.4 & 93.9 & 89.3 & 97.5 & 91.7 \\
        SCIDOCS & 92.9 & 93.0 & 89.7 & 97.6 & 92.6 \\
        SciFact & 96.9 & 97.6 & 96.3 & 97.9 & 96.7 \\
        \midrule
        \multicolumn{6}{l}{\textit{Reranking}} \\
        AskUbuntuDupQuestions & 99.2 & 102.2 & 96.6 & 97.6 & 99.1 \\
        SciDocsRR & 98.3 & 99.4 & 94.7 & 97.4 & 98.3 \\
        StackOverflowDupQuestions & 98.8 & 99.6 & 96.7 & 97.2 & 98.8 \\
        \midrule
        \multicolumn{6}{l}{\textit{STS}} \\
        BIOSSES & 98.8 & 105.1 & 94.0 & 97.6 & 99.2 \\
        SICK-R & 99.5 & 99.8 & 97.1 & 97.9 & 99.5 \\
        STS12 & 99.6 & 103.4 & 96.9 & 97.7 & 99.1 \\
        STS13 & 98.3 & 98.3 & 94.4 & 97.6 & 98.1 \\
        STS14 & 99.1 & 99.5 & 97.6 & 97.7 & 99.1 \\
        STS15 & 99.1 & 100.2 & 92.5 & 97.5 & 98.9 \\
        STS16 & 99.1 & 101.7 & 94.0 & 97.4 & 99.2 \\
        STSBenchmark & 99.2 & 100.4 & 97.6 & 97.8 & 99.1 \\
        \midrule
        \multicolumn{6}{l}{\textit{Pair Classification}} \\
        SprintDuplicateQuestions & 99.5 & 100.9 & 97.0 & 97.5 & 99.6 \\
        TwitterSemEval2015 & 99.3 & 101.7 & 96.4 & 97.8 & 99.5 \\
        TwitterURLCorpus & 99.4 & 100.0 & 98.3 & 97.8 & 99.4 \\
        \midrule
        \multicolumn{6}{l}{\textit{Summarization}} \\
        SummEval & 99.8 & 98.7 & 98.8 & 97.2 & 99.5 \\
        \midrule
        Mean & 97.5 & 99.9 & 90.6 & 97.6 & 97.6 \\
        \bottomrule
    \end{tabular}%
    }
\end{table*}

\begin{figure*}[t]
    \centering
    \includegraphics[width=0.95\linewidth]{figures/fig2_transfer_heatmap.pdf}
    \caption{Cross-task dimension importance transfer matrix for GTE-Large at 256 retained dimensions. Each cell shows retention when using task $i$'s ranking to prune for task $j$.}
    \label{fig:transfer_gte}
\end{figure*}

\section{Optimized--Random Gap Across the Number of Retained Dimensions}
\label{sec:gap_by_dim}

Table~\ref{tab:optimized_random_gap_by_dim} reports the optimized--random gap across retained dimensionalities. The gap is positive throughout the sweep, confirming that task-aware rankings identify more informative chunks than random selection. The gap widens as fewer dimensions are retained, indicating that performance becomes more sensitive to chunk choice when the selected subset is smaller.

This effect differs by model family. Retrieval-optimized Embedders and most Instruction-conditioned Embedders show modest gaps at 256 retained dimensions, but the gaps become larger at 64, 32, and 16 dimensions. RoBERTa-InBedder and General-purpose Language Model Backbones show larger gaps throughout the sweep, indicating that their task-relevant information is less evenly distributed across dimensions. Figure~\ref{fig:mechanism}(b) further shows that the normalized gap between optimized and worst-chunk selections shrinks quickly as more chunks are retained, meaning that once the retained subset is large enough, the exact chunk choice matters much less. 

\begin{table*}[!t]
    \centering
    \caption{Optimized--random gap across retained dimensions with chunk size $w=2$.}
    \label{tab:optimized_random_gap_by_dim}
    \small
    \resizebox{\textwidth}{!}{%
    \begin{tabular}{lcccccc}
        \toprule
        & \multicolumn{6}{c}{Retained dimensions} \\
        \cmidrule(lr){2-7}
        Model & $16$ & $32$ & $64$ & $128$ & $256$ & $512$ \\
        \midrule
        \multicolumn{7}{l}{\textit{Retrieval-optimized Embedders}} \\
        GTR-T5-Large      & +10.09\% & +9.69\% & +8.45\% & +6.88\% & +4.69\% & +2.35\% \\
        BGE-M3            & +11.06\% & +9.68\% & +8.10\% & +6.11\% & +4.61\% & +3.22\% \\
        GTE-Base          & +9.79\% & +8.44\% & +6.46\% & +4.73\% & +3.49\% & +1.82\% \\
        MxBai-Embed-Large & +9.66\% & +7.12\% & +5.11\% & +3.59\% & +2.59\% & +1.57\% \\
        GTE-Large         & +9.12\% & +7.31\% & +5.09\% & +3.31\% & +2.41\% & +1.72\% \\
        \midrule
        \multicolumn{7}{l}{\textit{Instruction-conditioned Embedders}} \\
        RoBERTa-InBedder  & +15.63\% & +14.33\% & +12.65\% & +10.92\% & +8.19\% & +6.18\% \\
        Qwen3-Embedding-0.6B   & +13.90\% & +11.99\% & +9.44\% & +6.99\% & +5.30\% & +3.10\% \\
        Stella EN 400M    & +9.82\% & +8.15\% & +6.09\% & +4.68\% & +3.22\% & +2.03\% \\
        Instructor-Large  & +9.20\% & +7.97\% & +6.05\% & +4.30\% & +2.89\% & +1.44\% \\
        \midrule
        \multicolumn{7}{l}{\textit{General-purpose Language Model Backbones}} \\
        RoBERTa-Large     & +20.33\% & +19.09\% & +18.45\% & +15.01\% & +11.75\% & +11.73\% \\
        BART-Base         & +19.22\% & +16.81\% & +14.90\% & +13.87\% & +10.08\% & +3.73\% \\
        \bottomrule
    \end{tabular}%
    }
\end{table*}

\section{Full Cross-Task Transfer Results Across Retained Dimensionalities}
\label{sec:multidim_appendix}

Table~\ref{tab:multidim_appendix} extends the main body results (Table~\ref{tab:cross_task_multidim}) to individual models. All values are mean retention across all source-target task pairs, including diagonal self-transfer.

\begin{table*}[t]
    \centering
    \caption{Cross-task transfer across six different retained dimensions for all 11 models with chunk size $w=2$. Each cell is mean retention over all source-target pairs, including diagonal self-transfer.}
    \label{tab:multidim_appendix}
    \small
    \resizebox{\textwidth}{!}{%
    \begin{tabular}{lcccccc}
        \toprule
        Model & Transfer@16 & Transfer@32 & Transfer@64 & Transfer@128 & Transfer@256 & Transfer@512 \\
        \midrule
        RoBERTa-Large & 67.0\% & 79.5\% & 93.5\% & 105.5\% & 114.0\% & 117.8\% \\
        RoBERTa-InBedder & 64.0\% & 75.9\% & 87.2\% & 96.2\% & 102.0\% & 104.8\% \\
        BART-Base & 64.4\% & 75.5\% & 85.5\% & 94.2\% & 100.1\% & 102.1\% \\
        GTE-Large & 67.8\% & 79.6\% & 88.8\% & 94.5\% & 97.6\% & 99.2\% \\
        Stella EN 400M & 67.5\% & 79.2\% & 88.5\% & 94.3\% & 97.5\% & 99.2\% \\
        MxBai-Embed-Large & 67.2\% & 78.8\% & 88.0\% & 93.9\% & 97.3\% & 99.1\% \\
        Instructor-Large & 64.4\% & 76.0\% & 86.1\% & 93.0\% & 97.0\% & 99.2\% \\
        Qwen3-Embedding-0.6B & 64.9\% & 76.6\% & 86.2\% & 92.9\% & 96.8\% & 99.1\% \\
        GTE-Base & 63.1\% & 75.0\% & 85.5\% & 92.7\% & 96.8\% & 99.1\% \\
        GTR-T5-Large & 62.2\% & 72.9\% & 83.2\% & 91.0\% & 96.1\% & 98.9\% \\
        BGE-M3 & 60.3\% & 70.7\% & 81.0\% & 89.2\% & 94.9\% & 98.3\% \\
        \bottomrule
    \end{tabular}%
    }
\end{table*}

\section{Chunk Size Sensitivity}
\label{sec:chunk_size}

We use chunk size $w = 2$ throughout the main experiments as a practical compromise between resolution and evaluation cost. For a typical 1024-dimensional model, evaluating single dimensions ($w = 1$) would require 1024 full task evaluations, while $w = 2$ cuts that to 512. Smaller chunks provide finer standalone scores but increase cost; larger chunks reduce cost but may hide within-chunk information. To verify that our conclusions are not sensitive to this choice, we sweep chunk sizes $w \in \{1, 2, 4, 8, 16, 32, 64\}$ on 5 models across 34 MTEB tasks.

\begin{table*}[h]
    \centering
    \caption{Optimized retention (\%) at dim=256 across chunk sizes averaged over 34 MTEB tasks. Random and Sequential baselines shown for reference. Bold marks the best chunk size per model.}
    \label{tab:chunk_sweep_full}
    \small
    \resizebox{\textwidth}{!} {
    \begin{tabular}{lccccccccc}
        \toprule
        & \multicolumn{7}{c}{Optimized retention at different chunk sizes $w$} & & \\
        \cmidrule(lr){2-8}
        Model & $w$=1 & $w$=2 & $w$=4 & $w$=8 & $w$=16 & $w$=32 & $w$=64 & Random & Seq \\
        \midrule
        \multicolumn{10}{l}{\textit{Retrieval-optimized Embedders}} \\
        GTE-Large & 99.9 & \textbf{99.9} & 99.6 & 99.4 & 99.4 & 99.2 & 98.9 & 97.5 & 97.5 \\
        GTE-Base & 99.9 & \textbf{100.1} & 99.8 & 99.1 & 98.8 & 98.3 & 98.2 & 96.6 & 96.6 \\
        GTR-T5-Large & \textbf{100.5} & 100.4 & 100.3 & 99.6 & 99.0 & 98.3 & 98.1 & 95.7 & 95.4 \\
        \midrule
        \multicolumn{10}{l}{\textit{Instruction-conditioned Embedders}} \\
        RoBERTa-InBedder & \textbf{108.7} & 108.5 & 107.4 & 107.0 & 108.1 & 108.0 & 107.4 & 100.3 & 101.3 \\
        Stella EN 400M & 100.1 & \textbf{100.3} & 100.2 & 100.2 & 99.8 & 99.2 & 99.0 & 97.1 & 97.1 \\
        \bottomrule
    \end{tabular}
    }
\end{table*}

Table~\ref{tab:chunk_sweep_full} shows the full results. Three findings emerge: (1) Optimized retention is nearly identical at $w = 1$ and $w = 2$ for all 5 models, indicating that the main optimized–random gaps are not an artifact of using paired dimensions. (2)~The optimized--random gap degrades gracefully as chunk size increases, and even $w = 64$ retains a positive gap for all models. (3)~The association between the optimized--random gap and training paradigm persists across chunk sizes. The gaps remain small for Retrieval-optimized Embedders and Stella EN 400M ($+1.4$--$4.8$\%), but are larger for RoBERTa-InBedder ($+6.7$--$8.4$\%). Thus, varying the chunk size does not change the main empirical conclusion: the value of task-aware pruning remains associated with training paradigm. These results support our choice of $w = 2$ for the main experiments.

\section{Random Variance and Tail Risk}
\label{sec:random_variance}

We quantify the full distribution of random selection performance using 100 seeds per combination across 9 models and 5 retained dimensionalities ($d \in \{16, 32, 64, 128, 256\}$).

\begin{table}[!h]
    \centering
    \caption{Random selection distribution statistics across 100 seeds (mean over 20 tasks per model). CV = coefficient of variation. P5 = 5th percentile. CVaR-5 = conditional value at risk at 5\%.}
    \label{tab:random_variance}
    \small
    \resizebox{\columnwidth}{!}{%
    \begin{tabular}{lccccc}
        \toprule
        Model & Dim & Mean $\pm$ Std & CV & P5 & CVaR-5 \\
        \midrule
        \multicolumn{6}{l}{\textit{Retrieval-optimized Embedders}} \\
        GTE-Large & 256 & (97.2 $\pm$ 1.3)\% & 1.3\% & 95.1\% & 95.0\% \\
        GTE-Large & 64 & (86.9 $\pm$ 2.1)\% & 2.6\% & 82.9\% & 82.7\% \\
        GTE-Large & 16  & (63.9 $\pm$ 3.1)\% & 6.8\% & 58.7\% & 58.4\% \\
        \midrule
        \multicolumn{6}{l}{\textit{Instruction-conditioned Embedders}} \\
        Stella EN 400M & 256 & (96.6 $\pm$ 1.2)\% & 1.2\% & 94.6\% & 94.5\% \\
        Stella EN 400M & 64 & (87.3 $\pm$ 2.3)\% & 2.8\% & 83.7\% & 83.5\% \\
        Stella EN 400M & 16 & (63.7 $\pm$ 3.4)\% & 7.9\% & 58.2\% & 58.0\% \\
        \midrule
        \multicolumn{6}{l}{\textit{General-purpose Language Model Backbones}} \\
        RoBERTa-Large & 256 & (114.3 $\pm$ 8.0)\% & 5.7\% & 97.0\% & 96.4\% \\
        RoBERTa-Large & 64 & (99.4 $\pm$ 8.7)\% & 8.4\% & 83.2\% & 82.5\% \\
        RoBERTa-Large & 16 & (76.2 $\pm$ 7.7)\% & 15.1\% & 62.4\% & 61.7\% \\
        BART-Base & 256 & (97.8 $\pm$ 6.6)\% & 6.9\% & 86.2\% & 85.7\% \\
        BART-Base & 64  & (82.9 $\pm$ 8.0)\% & 11.5\% & 70.0\% & 69.3\% \\
        BART-Base & 16 & (62.1 $\pm$ 7.0)\% & 18.5\% & 50.0\% & 49.4\% \\
        \bottomrule
    \end{tabular}%
    }
\end{table}

Table~\ref{tab:random_variance} summarizes the representative results. At 256 retained dimensions, GTE-Large and Stella EN 400M are stable, with CV around 1.2--1.3\% and P5 about 2 percentage points below the mean. Seed sensitivity increases at dim=16, but the lower tail remains moderate for these models. In contrast, General-purpose Language Model Backbones show much larger variance: at 256 retained dimensions, CV reaches 5.7--6.9\%, and at dim=16 it rises to 15.1--18.5\%. Thus, random selection has low tail risk for embedding-specific models at moderate compression, while tail risk becomes more severe for General-purpose Language Model Backbones and aggressive compression.

\begin{figure*}[h]
    \centering
    \includegraphics[width=0.95\linewidth]{figures/fig_random_variance.pdf}
    \caption{Retention distributions under random dimension selection across tasks. Each box plot shows the task-level spread of mean retention for retained dimensions $\in$ \{64, 128, 256\}. Triangle markers ($\triangle$) indicate P5 retention; diamond markers ($\diamond$) indicate CVaR-5 retention. Tail risk increases sharply as the retained dimension decreases.}
    \label{fig:random_variance_app}
\end{figure*}

\section{Task-Level Low-Retention Outcomes}
\label{sec:task_low_retention}

Retention ratio is useful for comparing tasks with different raw metrics, but ratio averaging may hide failures. Therefore, Table~\ref{tab:low_retention_random} reports low-retention summaries at 256 retained dimensions. P25, P10, and P5 describe progressively lower parts of the task distribution, while the worst-task column identifies where each model is most sensitive.


\begin{table*}[t]
    \centering
    \caption{Random low-retention summaries at 256 retained dimensions. P25, P10, and P5 are task-level percentiles. Worst task reports the lowest-retention task and its retention.}
    \label{tab:low_retention_random}
    \scriptsize
    \resizebox{\textwidth}{!}{%
    \begin{tabular}{lrrrrlrr}
        \toprule
        Model & Mean & P25 & P10 & P5 & Worst task & $<90$ & $>100$ \\
        \midrule
        \multicolumn{8}{l}{\textit{Retrieval-optimized Embedders}} \\
        GTR-T5-Large & 95.7\% & 92.6\% & 90.3\% & 89.3\% & SCIDOCS (88.4\%) & 3 & 1 \\
        BGE-M3 & 94.6\% & 91.9\% & 86.5\% & 84.0\% & SCIDOCS (80.2\%) & 7 & 0 \\
        GTE-Base & 96.6\% & 94.4\% & 93.1\% & 92.5\% & MTOPIntentClassification (91.2\%) & 0 & 0 \\
        MxBai-Embed-Large & 97.2\% & 96.1\% & 93.4\% & 92.7\% & NFCorpus (91.1\%) & 0 & 0 \\
        GTE-Large & 97.5\% & 96.8\% & 94.0\% & 92.4\% & TwentyNewsgroupsClustering (90.8\%) & 0 & 0 \\
        \midrule
        \multicolumn{8}{l}{\textit{Instruction-conditioned Embedders}} \\
        RoBERTa-InBedder & 100.3\% & 96.3\% & 94.2\% & 93.3\% & TwentyNewsgroupsClustering (91.4\%) & 0 & 11 \\
        Qwen3-Embedding-0.6B & 95.5\% & 92.6\% & 90.6\% & 88.8\% & NFCorpus (83.6\%) & 3 & 0 \\
        Stella EN 400M & 97.1\% & 95.7\% & 93.0\% & 92.2\% & NFCorpus (89.8\%) & 1 & 1 \\
        Instructor-Large & 97.0\% & 95.8\% & 93.6\% & 91.9\% & NFCorpus (89.0\%) & 1 & 1 \\
        \midrule
        \multicolumn{8}{l}{\textit{General-purpose Language Model Backbones}} \\
        RoBERTa-Large & 111.5\% & 96.1\% & 90.8\% & 89.9\% & Banking77Classification (87.2\%) & 2 & 18 \\
        BART-Base & 97.9\% & 95.1\% & 93.0\% & 91.7\% & CQADupstackEnglishRetrieval (86.8\%) & 1 & 6 \\
        \bottomrule
    \end{tabular}%
    }
\end{table*}

At 256 retained dimensions, 10 of 11 models have P10 retention above 90\%, and 7 of 11 models have at most one task below 90\%. This confirms that the mean retention is not driven only by a few extremely easy or difficult tasks: most of the task distribution remains close to full-dimensional performance.

The lower-retention cases are also structured rather than arbitrary. Across all 11 models, the task-model pairs below 90\% come mostly from retrieval tasks (9 pairs) and classification tasks (7 pairs), with only two clustering pairs and none from other task categories. Worst-task entries show the same pattern: NFCorpus (retrieval) is the lowest-retention task for four models, SCIDOCS (retrieval) for two models, and TwentyNewsgroupsClustering (clustering) for two models. Thus, the low end of the distribution points to a small set of sensitive benchmark tasks rather than a collapse of selection strategy.

\section{Analysis of Alternative Importance Definitions}
\label{sec:loo}

We compare the standalone chunk score used in the main analysis with three alternative importance definitions. As in the main text, $S^{M,T}(A)$ denotes the task score of model $M$ on task $T$ when retaining dimension set $A$, and $C_i$ denotes the $i$-th chunk, $\mathrm{Full}$ denotes the full set of embedding dimensions. We define $\mathcal{A}_{-i}$ as the collection of retained dimension sets that are unions of chunks and do not contain $C_i$.

\paragraph{Leave-One-Out importance.}
Leave-One-Out (LOO) importance measures the performance drop caused by removing chunk $C_i$ from the full representation:
\begin{equation}
\ell_i^{M,T}
=
S^{M,T}(\mathrm{Full})
-
S^{M,T}(\mathrm{Full}\setminus C_i)
\label{eq:loo_importance}
\end{equation}
Thus, LOO evaluates whether $C_i$ is necessary in the presence of all other chunks.

\paragraph{Marginal importance.}
Marginal importance measures the expected gain from adding chunk $C_i$ to a retained dimension set that does not already contain it:
\begin{equation}
m_i^{M,T}
=
\mathbb{E}_{A\sim \mathcal{A}_{-i}}
\left[
S^{M,T}(A\cup C_i)
-
S^{M,T}(A)
\right]
\label{eq:marginal_importance}
\end{equation}
This captures whether $C_i$ provides additional value when combined with other retained dimensions.

\paragraph{Shapley importance.}
Shapley importance averages the marginal contribution of $C_i$ over all possible retained dimension sets that exclude it:
\begin{equation}
\phi_i^{M,T}
=
\sum_{A\in \mathcal{A}_{-i}}
w(A)
\left[
S^{M,T}(A\cup C_i)-S^{M,T}(A)
\right]
\label{eq:shapley_importance}
\end{equation}
where $w(A)=|A|_c!(N-|A|_c-1)!/N!$ is the standard Shapley weight and $|A|_c$ denotes the number of chunks contained in $A$. This gives a context-averaged attribution of chunk importance and accounts for interactions among chunks. In practice, marginal and Shapley importance are estimated by sampling retained dimension sets from $\mathcal{A}_{-i}$.

Table~\ref{tab:importance_agreement} reports model-level Spearman correlation between standalone and alternative importance definitions and mean entropy per definition. Table~\ref{tab:loo_detailed} further shows per-task results on RoBERTa-Large. The weak correlation between standalone and Leave-One-Out/marginal/Shapley importance suggests that chunk importance is context-dependent. Standalone entropy is near-maximum (0.985--1.000), while Shapley entropy is substantially lower (0.655--0.933), revealing that interaction-aware importance detects much more structure. 


\begin{table*}[h]
    \centering
    \caption{Spearman correlation between standalone and alternative importance definitions at 256 retained dimensions (chunk size $w=4$), and mean entropy per definition.}
    \label{tab:importance_agreement}
    \small
    \resizebox{\textwidth}{!}{%
    \begin{tabular}{lccccccc}
        \toprule
        & \multicolumn{3}{c}{Correlation with Standalone} & \multicolumn{4}{c}{Entropy} \\
        \cmidrule(lr){2-4} \cmidrule(lr){5-8}
        Model & LOO & Marginal & Shapley & Standalone & LOO & Marginal & Shapley \\
        \midrule
        GTE-Large & 0.094 & 0.155 & 0.132 & 0.999 & 0.936 & 0.939 & 0.625 \\
        BGE-M3 & 0.175 & 0.251 & 0.181 & 0.998 & 0.932 & 0.939 & 0.597 \\
        Stella EN 400M & 0.124 & 0.209 & 0.157 & 0.999 & 0.918 & 0.944 & 0.610 \\
        RoBERTa-Large & 0.091 & 0.296 & 0.235 & 0.997 & 0.838 & 0.891 & 0.660 \\
        RoBERTa-InBedder & 0.096 & 0.248 & 0.179 & 0.997 & 0.888 & 0.927 & 0.654 \\
        \bottomrule
    \end{tabular}%
    }
\end{table*}

\begin{table*}[h]
    \centering
    \caption{Per-task importance definition comparison for RoBERTa-Large at 256 retained dimensions (chunk size $w=4$). Spearman correlations and normalized entropy for standalone, Leave-one-out, marginal, and Shapley importance. Low standalone-to-alternative correlations confirm importance is context-dependent.}
    \label{tab:loo_detailed}
    \small
    \resizebox{\textwidth}{!}{%
    \begin{tabular}{lcccccccc}
        \toprule
        & \multicolumn{4}{c}{Spearman $\rho$ with Standalone} & \multicolumn{4}{c}{Normalized Entropy} \\
        \cmidrule(lr){2-5} \cmidrule(lr){6-9}
        Task & LOO & Marginal & Shapley & $p$ (best) & Standalone & LOO & Marginal & Shapley \\
        \midrule
        ImdbClassification & $-$0.080 & 0.252 & 0.176 & $<$0.001 & 1.000 & 0.943 & 0.942 & 0.707 \\
        Banking77Classification & $-$0.001 & 0.078 & 0.156 & 0.012 & 0.996 & 0.949 & 0.914 & 0.860 \\
        TwentyNewsgroupsClustering & $-$0.038 & 0.137 & 0.132 & 0.029 & 0.998 & 0.922 & 0.950 & 0.932 \\
        MedrxivClusteringS2S & 0.018 & 0.039 & 0.073 & 0.247 & 0.999 & 0.952 & 0.945 & 0.842 \\
        STSBenchmark & 0.387 & 0.704 & 0.412 & $<$0.001 & 0.998 & 0.343 & 0.677 & 0.655 \\
        BIOSSES & 0.260 & 0.568 & 0.459 & $<$0.001 & 0.985 & 0.921 & 0.920 & 0.783 \\
        \bottomrule
    \end{tabular}%
    }
\end{table*}

\paragraph{Why standalone scoring?}
We use standalone chunk scores in the main text because our goal is to diagnose whether individual chunks can carry task-relevant information. In highly redundant embeddings, evaluating the effect of adding or removing one chunk from a large retained set often produces only minor changes in MTEB metrics, since many other chunks can compensate for the removed information. Such low-contrast marginal scores can make chunk rankings unstable. Standalone evaluation removes this compensation effect by asking how much task-relevant signal each chunk can support on its own, yielding a more discriminative diagnostic of isolated chunk informativeness.

\section{Practical Selection Baselines}
\label{sec:practical_baselines}

The main experiments use optimized selection as a diagnostic, but it requires evaluating every chunk on each target task. We therefore ask whether lower-cost selection rules can approximate the same task-relevant dimension structure. As in the main text, $C_i$ denotes the i-th contiguous chunk of size $w$, and let $m=k/w$ be the number of retained chunks. For a baseline that assigns each chunk a score $u_i$, we retain
\begin{equation}
A_{\mathrm{base}}(k)
=
\operatorname{Top}_{m}\left(\{u_i\}_{i=1}^{N}\right)
\end{equation}
where the selected chunks correspond to the retained dimensions. 

\paragraph{Sequential selection.}
Sequential selection retains the first $m$ chunks in the model's original dimension order:
\begin{equation}
A_{\mathrm{seq}}(k)
=
\{C_1,\ldots,C_m\}
\end{equation}
This tests whether the native dimension order already encodes a useful importance ranking, as in Matryoshka-style representations. 

\paragraph{Magnitude-based selection.}
Magnitude-based selection scores each dimension by the $\ell_2$ norm of its corresponding column in the embedding weight matrix. Let $\mathbf{W}\in\mathbb{R}^{V\times D}$ be the embedding matrix. The dimension score is:
\begin{equation}
u_d^{\mathrm{mag}}
=
\|\mathbf{W}_{:,d}\|_2
\end{equation}

This baseline tests whether dimensions with larger embedding-weight magnitude are also more important for downstream task performance.

\paragraph{Gradient saliency.}
Gradient saliency uses local task-loss sensitivity to score dimensions. For task $T$, embedding $\mathbf{e}(x)$, and task loss $\mathcal{L}_T$, we define
\begin{equation}
u_d^{\mathrm{grad}}
=
\mathbb{E}_{x}
\left[
\left|
e_d(x)
\frac{\partial \mathcal{L}_T}{\partial e_d(x)}
\right|
\right]
\end{equation}
The corresponding chunk score is
\begin{equation}
u_i^{\mathrm{grad}}
=
\sum_{d\in C_i}u_d^{\mathrm{grad}} 
\end{equation}
This method tests whether dimensions with larger local influence on the task loss also preserve more task-relevant information after selection.

\paragraph{Activation variance.}
Activation variance scores dimensions by how much their embedding values vary across input examples:
\begin{equation}
u_d^{\mathrm{var}}
=
\operatorname{Var}_{x}\left[e_d(x)\right]
\end{equation}
The corresponding chunk score is
\begin{equation}
u_i^{\mathrm{var}}
=
\sum_{d\in C_i}u_d^{\mathrm{var}} 
\end{equation}
The intuition is that dimensions with larger variation across examples may carry more discriminative information.

\paragraph{Learned mask selection.}
The learned-mask baseline directly optimizes a continuous mask over embedding dimensions. Let $\mathbf{m}\in[0,1]^D$ be applied to the embedding:
\begin{equation}
\tilde{\mathbf{e}}(x)
=
\mathbf{m}\odot\mathbf{e}(x)
\end{equation}
We optimize the mask with a task loss and an $\ell_1$ sparsity penalty:
\begin{equation}
\min_{\mathbf{m}\in[0,1]^D}
\;
\mathbb{E}_{(x,y)}
\left[
\mathcal{L}_T(\mathbf{m}\odot\mathbf{e}(x),y)
\right]
+
\lambda\|\mathbf{m}\|_1 
\end{equation}
After optimization, dimensions are ranked by their mask values.

\paragraph{Projection-based compression.}
PCA and random projection are included as projection-based compression baselines, although they do not retain a subset of original dimensions. For PCA, given centered embeddings $\mathbf{e}-\boldsymbol{\mu}$ and the top $k$ principal directions $\mathbf{U}_k\in\mathbb{R}^{D\times k}$, we use
\begin{equation}
\mathbf{z}^{\mathrm{PCA}}
=
\mathbf{U}_k^\top(\mathbf{e}-\boldsymbol{\mu})
\end{equation}
For random projection, with projection matrix $\mathbf{G}\in\mathbb{R}^{D\times k}$, we use
\begin{equation}
\mathbf{z}^{\mathrm{RP}}
=
\mathbf{G}^\top\mathbf{e}
\end{equation}
We evaluate these compressed representations using the same retention ratio as in the main text, replacing the retained-dimension representation with the projected representation.

\begin{table*}[!t]
    \centering
    \caption{Chunk importance distribution across five representative models. Optimized--random gap is measured at 256 retained dimensions, with chunk size $w=2$.}
    \label{tab:inbedder_mechanism}
    \small
    \resizebox{\textwidth}{!}{%
    \begin{tabular}{lccccc}
        \toprule
        Metric & GTE-Large & Stella EN 400M & RoBERTa-InBedder & BART-Base & RoBERTa-Large \\
        \midrule
        Normalized entropy & 0.994 & 0.993 & 0.990 & 0.989 & 0.989 \\
        Gini coefficient & 0.101 & 0.107 & 0.124 & 0.130 & 0.131 \\
        Chunk importance CV & 0.186 & 0.195 & 0.242 & 0.252 & 0.255 \\
        Top-50\% concentration & 0.573 & 0.576 & 0.593 & 0.598 & 0.598 \\
        \midrule
        Optimized--random gap & +2.41\% & +3.22\% & +8.19\% & +10.08\% & +11.75\% \\
        \bottomrule
    \end{tabular}%
    }
\end{table*}


\paragraph{Results.}
\begin{table}[t]
    \centering
    \caption{Mean retention of gradient saliency, activation variance, and learned-mask selection on 8 representative tasks at 256 dimensions. Gradient and activation variance can be actively harmful for RoBERTa-InBedder and RoBERTa-Large.}
    \label{tab:strong_baselines}
    \small
    \resizebox{\columnwidth}{!}{%
    \begin{tabular}{lcccc}
        \toprule
        Model & Random & Gradient & Activation Variance & Mask \\
        \midrule
        \multicolumn{5}{l}{\textit{Retrieval-optimized Embedders}} \\
        GTE-Large & 99.5\% & 99.5\% & 99.7\% & 98.9\% \\
        BGE-M3 & 96.8\% & 97.1\% & 97.6\% & 96.4\% \\
        \midrule
        \multicolumn{5}{l}{\textit{Instruction-conditioned Embedders}} \\
        Stella EN 400M & 98.8\% & 99.3\% & 98.4\% & 98.3\% \\
        RoBERTa-InBedder & 100.1\% & 94.8\% & 93.3\% & 98.5\% \\
        \midrule
        \multicolumn{5}{l}{\textit{General-purpose Language Model Backbones}} \\
        RoBERTa-Large & 102.9\% & 93.4\% & 92.4\% & 102.0\% \\
        \bottomrule
    \end{tabular}%
    }
\end{table}

\begin{table}[t]
    \centering
    \caption{Random selection vs. projection-based compression, comparing random selection (RS), PCA, and random projection (RP) at 64 and 256 dimensions across representative models.}
    \label{tab:pca_comparison}
    \small
    \resizebox{\columnwidth}{!}{%
    \begin{tabular}{llccc}
        \toprule
        Model & Dim & PCA & RP & RS \\
        \midrule
        \multirow{2}{*}{GTE-Large}
            & 64  & 90.2\% & 92.9\% & 93.1\% \\
            & 256 & 98.0\% & 98.7\% & 99.0\% \\
        \multirow{2}{*}{BGE-M3}
            & 64  & 85.9\% & 87.2\% & 87.8\% \\
            & 256 & 95.8\% & 96.4\% & 96.1\% \\
        \multirow{2}{*}{Stella EN 400M}
            & 64  & 92.0\% & 93.9\% & 94.4\% \\
            & 256 & 98.5\% & 98.6\% & 98.9\% \\
        \multirow{2}{*}{RoBERTa-InBedder}
            & 64  & 80.8\% & 90.4\% & 96.7\% \\
            & 256 & 95.6\% & 98.4\% & 98.0\% \\
        \multirow{2}{*}{RoBERTa-Large}
            & 64  & 80.6\% & 91.2\% & 98.4\% \\
            & 256 & 99.9\% & 99.6\% & 100.2\% \\
        \bottomrule
    \end{tabular}%
    }
\end{table}

\begin{figure*}[t]
    \centering
    \includegraphics[width=0.95\linewidth]{figures/fig9_magnitude_comparison.pdf}
    \caption{(a) Spearman rank correlation between magnitude-based dimension ranking and task-specific importance ranking across five models. All mean correlations remain near zero. (b--c) Per-task magnitude vs.\ random retention. Points below the diagonal indicate tasks where random outperforms magnitude.}
    \label{fig:magnitude}
\end{figure*}

\begin{figure*}[t]
    \centering
    \includegraphics[width=0.95\linewidth]{figures/fig10_all_methods_comparison.pdf}
    \caption{Comparison of five pruning methods at 256 retained dimensions for three representative models. Error bars show standard deviations across tasks. Random, Sequential, and Magnitude cluster together, while Anti-optimized drops to 86--89\%. The optimized selection strategy lies only slightly above the cluster for GTE-Large and Stella EN 400M, but noticeably farther above it for RoBERTa-InBedder. Magnitude slightly underperforms Random on GTE-Large and Stella EN 400M, but shows a much larger drop on RoBERTa-InBedder. }
    \label{fig:all_methods}
\end{figure*}

Table~\ref{tab:strong_baselines} reports the representative-task results for gradient saliency, activation variance, and learned-mask selection, while Table~\ref{tab:pca_comparison} compares PCA, random projection, and random selection. Sequential and magnitude-based results are reported in Figure~\ref{fig:magnitude} and~\ref{fig:all_methods}.

Across these baselines, simple task-agnostic signals (prefix order, weight magnitude, activation variance, and random projections) and lightweight task-aware signals (gradient saliency and learned masks) do not consistently recover a substantially better retained subset than random selection at moderate compression. This supports the main diagnostic conclusion: for many embedding-specific models, task-relevant information is distributed broadly enough that many retained subsets perform similarly.

\begin{figure*}[t]
    \centering
    \includegraphics[width=0.95\linewidth]{figures/fig7_redundancy_mechanism.pdf}
    \caption{Representative views of the redundancy mechanism for the three representative models. (a) Histograms of per-task normalized entropy. (b) Normalized optimized--worst selection gap as the number of selected chunks increases.}
    \label{fig:mechanism}
\end{figure*}

\begin{figure*}[t]
    \centering
    \includegraphics[trim=0 26cm 0 1cm, clip, width=0.95\linewidth]{figures/fig_basis_independence.pdf}
    \caption{Basis-dependence analysis of random selection. We apply six basis transformations to the embedding space (original, three random orthogonal rotations, PCA, and whitening) and measure the sequential--random gap and cosine--correlation preservation after pruning. In every basis, the gap stays below 1\% and cosine correlation exceeds 0.94, suggesting that the broad distribution of task-relevant information is not merely an artifact of one coordinate system.}
    \label{fig:basis_independence}
\end{figure*}


\section{Retrieval System Cost}
\label{sec:retrieval_cost}
Using FAISS FlatIP, Random-256 retains 87.5--98.4\% of full-dim nDCG while reducing raw vector storage by 4$\times$ (Table~\ref{tab:retrieval_cost}). Random projection is included as a projection-based comparator, but it requires applying a projection matrix. For GTE-Large, Random-256 reduces P50 latency from 0.16--0.23ms to 0.04--0.08ms, a 2.9--4.3$\times$ speedup.

\begin{table}[t]
    \centering
    \caption{Retrieval cost--quality tradeoff at dim=256 for GTE-Large using FAISS FlatIP. nDCG@10 and Recall report retention relative to the full-dimensional baseline for each task.}
    \label{tab:retrieval_cost}
    \resizebox{\columnwidth}{!}{%
    \begin{tabular}{llcccc}
        \toprule
        Task & Method & nDCG@10 & Mem (MB) & Latency (ms) & Recall \\
        \midrule
        \multirow{3}{*}{NFCorpus} & Full-1024 & 100.0\% & 14.2 & 0.16 & 100.0\% \\
        & Random-256 & 98.4\% & 3.5 & 0.04 & 96.7\% \\
        & RP-256 & 85.3\% & 3.5 & 0.38 & 82.9\% \\
        \midrule
        \multirow{3}{*}{ArguAna} & Full-1024 & 100.0\% & 19.5 & 0.22 & 100.0\% \\
        & Random-256 & 93.7\% & 4.9 & 0.05 & 95.9\% \\
        & RP-256 & 94.0\% & 4.9 & 0.40 & 93.8\% \\
        \midrule
        \multirow{3}{*}{SciFact} & Full-1024 & 100.0\% & 19.5 & 0.23 & 100.0\% \\
        & Random-256 & 87.5\% & 4.9 & 0.08 & 93.8\% \\
        & RP-256 & 83.4\% & 4.9 & 0.80 & 88.1\% \\
        \bottomrule
    \end{tabular}%
    }
\end{table}


\section{Short-Answer-Conditioned Supervision Is Associated with Greater Dimension Structure}
\label{sec:Inbedder-mechanism}
\paragraph{Characteristics of the exception.}
To understand why RoBERTa-InBedder differs, we compare its dimension importance distribution with one Retrieval-Optimized Embedder (GTE-Large) and one Instruction-conditioned Embedder (Stella EN 400M) (Table~\ref{tab:inbedder_mechanism}, Figure~\ref{fig:mechanism}). Its Gini coefficient is 0.124 vs.\ 0.101--0.107 for GTE-Large and Stella EN 400M, its chunk-importance coefficient of variation is 0.242 vs.\ 0.186--0.195. Its normalized entropy is 0.990 vs.\ 0.993--0.994. While these absolute differences are small, the relative differences are consistent: RoBERTa-InBedder's importance distribution is measurably less uniform. Table~\ref{tab:inbedder_mechanism} also demonstrates that General-purpose Language Model Backbones show even lower entropy and higher Gini, confirming that training paradigm influences dimension structure.

\paragraph{Mechanism: short-answer-conditioned supervision reduces uniformity.}
RoBERTa-InBedder shows consistently higher concentration of importance across all metrics than Retrieval-optimized Embedders and other Instruction-conditioned Embedders. A plausible interpretation is that its supervision differs from standard embedding training in both objective and signal breadth. InBedder instantiates instruction-following embedding by treating the instruction as a question and learning representations through expected answers; in practice, it is fine-tuned on abstractive QA data, with answers simplified to an average length of only 2.89 tokens \citep{inbedder}. Such short-answer-conditioned supervision may provide a narrower training signal than contrastive or multi-task embedding objectives, causing the model to align a smaller subset of dimensions with latent answer features while leaving other dimensions less consistently useful. This interpretation is consistent with prior evidence that fine-tuning can substantially reshape embedding geometry rather than simply improving isotropy \citep{rajaee2021fine}.

\section{Theoretical perspective: near-uniform importance as a training-induced pattern}
\label{theoretical_perspective}
Under a contrastive loss $\mathcal{L} = -\log[\exp(\mathbf{z}_i \cdot \mathbf{z}_j) / \sum_{k} \exp(\mathbf{z}_i \cdot \mathbf{z}_k)]$ where $\mathbf{z}_i$ is the representation of anchor $i$, $\mathbf{z}_j$ is the positive, and $\mathbf{z}_k$ ranges over all positives and negatives in the batch, each dimension $d$ contributes additively to the dot product. When the training data is broad and diverse, no single dimension can specialize too strongly without hurting performance elsewhere, so the contributions tend to balance out. The RoBERTa-InBedder exception fits this picture in that its non-contrastive, answer-conditioned fine-tuning on very short QA answers is associated with a less even importance distribution and a larger task-aware advantage. Taken together, this reasoning suggests a clear prediction for future work: models trained with less diverse supervision should display less uniform importance distributions and, as a result, larger optimized–random gaps. A rigorous test of this prediction would require controlling for confounding factors by holding architecture and data fixed while varying the training objective, or by comparing matched model pairs that share the same embedding head and normalization. We leave this controlled investigation to future work.

We also repeat the pruning analysis after several basis transformations, including random orthogonal rotations, PCA, and whitening. The sequential--random gap stays below 1\% in every basis (Figure~\ref{fig:basis_independence}). This suggests that the observed robustness of random selection is not merely a quirk of the native coordinate basis.


\end{document}
